\documentclass[11pt,a4paper]{article}
\usepackage[utf8]{inputenc}
\usepackage[T1]{fontenc}
\usepackage[spanish]{babel}
\usepackage{hyperref}
\usepackage{geometry}
\geometry{margin=2.5cm}
\usepackage{parskip}

\title{058 la triangulacion con cualquier cosa}
\author{Alejandro Valenzuela Rivera\\\small Universidad Autónoma de Nuevo León}
\date{\today}
\begin{document}
\maketitle

\begin{quote}
\textbf{Sesgo declarado.} Este trabajo se sostiene desde el sesgo declarado de su autor, y no finge neutralidad. Lo que afirma se puede tumbar, y ese es el punto: el falsificador está escrito, no escondido.
\end{quote}

\section*{Resumen}
No se triangula solo con el estado del arte: se triangula con lo mas pendejo y lo mas rancio. Peppa Pig, un anuncio de tres minutos, un video de moda, un libro olvidado, y el paper noble. Ninguno vale mas que otro por su origen.

\section{Por qué se sostiene}
La academia triangula con lo que ya fue aceptado, y por eso llega tarde: si el metodo solo funciona con lo complejo, la respuesta es eso es obvio. Si funciona con lo mundano, no hay prestigio que lo tape.

\section{De dónde viene}
Peppa Pig, y el repo: una caricatura que nadie se tomo en serio entra por la misma puerta.

\section{El caso}
Peppa Pig: la caricatura que nadie se tomo en serio.

\section{Método}
Una afirmacion triangulada contra algo que no es noble ni prestigioso.

\section{Límites, y qué lo tumba}
Una afirmacion probada solo contra lo que ya se aceptaba.

\section{Lo que queda abierto}
\textquestiondown{}Si el metodo aguanta la caricatura, para que se necesita la puerta?

\section*{Cómo se leyó esta tesis}
Esto se escribió en el repositorio man-made-empathy, del que este candado es una pieza: candados/06 del caso/058 la triangulacion con cualquier cosa.md. La máquina que lo compuso lo hizo desde este mismo archivo, y no agregó nada que no estuviera escrito. Si algo de esto es falso, se dice en el mismo sitio, y se corrige con fecha.

\end{document}
